% !TeX root = ../../Book3.tex
% 纲要标题：### D.4 形式纤维、卓越性与逼近
% 纲要位置：工作记录/计划纲要.md:2936

\section{形式纤维、卓越性与逼近}
\label{b3:sec:D04}

一组相容的形式近似不一定就是原环中的精确解，完备化的平坦性也不保证其各条纤维正则。本节在明确的底环假设下讨论这些差别，并区分有限阶逼近与实际收敛。

% 纲要内容（仅作写作计划，尚非教材正文）：
% * 完备化映射的形式纤维与卓越性的用途
% * 卓越性的来源与稳定性采用准确引用；Artin逼近限定域上有限型局部代数的Hensel化，保留Jacobian修正的可核验末步
% * 逼近意味着模高次极大理想的一致，不是任意形式解本身收敛或代数化
% * 与附录 C 的奇点理论结合阅读，不把全部消解理论作为代数学共同基础
% #### D.4.1 完备化的形式纤维与卓越性
% #### D.4.2 Artin 逼近：域上有限型情形
% #### D.4.3 有限阶逼近、收敛性与代数化
% #### D.4.4 奇点与消解理论导读


\subsection{完备化的形式纤维与卓越性}
\label{b3:subsec:D04:01}

忠实平坦性保证完备化能够检测理想与有限生成模的关系，但并不说明它的每条纤维都正则。后一个问题涉及完备化在不同素点上新增加的元素。只检查闭点的剩余域，不能回答这个问题。

\begin{definition}\label{b3:def:formal-fiber}
设 \((A,\mathfrak m)\) 为 Noether 局部环，\(\widehat A\) 为其 \(\mathfrak m\)-进完备化。对 \(\mathfrak p\in\operatorname{Spec}A\)，环
\[
\widehat A\otimes_A\kappa(\mathfrak p)
\]
称为 \(A\to\widehat A\) 在 \(\mathfrak p\) 处的\term[formalfiber]{形式纤维}[formal fiber]，其中
\(\kappa(\mathfrak p)=\operatorname{Frac}(A/\mathfrak p)\) 是该素点的剩余域。
\end{definition}

闭点处的形式纤维总是
\(\widehat A/\mathfrak m\widehat A\simeq A/\mathfrak m\)；其他素点的形式纤维不能由这一个域代替。例如
\(A=k[x]_{(x)}\) 在零素理想处的形式纤维为 \(k((x))\)，这是 \(k(x)\)-代数；附录 C 的\secref{b3:sec:C01}逐项计算了这两条纤维。若 \(A\) 本身已完备，则完备化映射为恒等，每条形式纤维才都只是相应的 \(\kappa(\mathfrak p)\)。

\begin{definition}\label{b3:def:excellent-ring}
设 \(A\) 为交换 Noether 环。
\begin{enumerate}
\item 若对每对素理想 \(\mathfrak p\subseteq\mathfrak q\)，连接二者的饱和素理想链都有相同长度，则称 \(A\) 为悬链环；饱和表示相邻两项之间不能再插入素理想。若每个有限型 \(A\)-代数都为悬链环，则称 \(A\) 具有普适链性质。
\item 设 \(F\) 为域。一个 Noether \(F\)-代数 \(B\)，若对每个有限域扩张 \(F'/F\)，环 \(B\otimes_FF'\) 都正则，则称 \(B\) 在 \(F\) 上几何正则。
\item 若对每个 \(\mathfrak p\in\operatorname{Spec}A\)，映射
\[
A_{\mathfrak p}\longrightarrow\widehat{A_{\mathfrak p}}
\]
的每条形式纤维在相应剩余域上都几何正则，则称 \(A\) 为 \(G\) 环。
\item 若对每个有限型 \(A\)-代数 \(B\)，集合
\(\operatorname{Reg}(B)=\{\mathfrak q:B_{\mathfrak q}\text{ 正则}\}\)
在 \(\operatorname{Spec}B\) 中开，则称 \(A\) 满足 \(J\)-2 条件。
\end{enumerate}
若 \(A\) 同时具有普适链性质、\(G\) 环性质及 \(J\)-2 条件，则称它为\term[excellentring]{卓越环}[excellent ring]。
\end{definition}

这里悬链环对应 catenary ring，只表示素链长度条件，不表示所有理想线性排序。几何正则要求允许有限的不可分域扩张，不能只检查可分扩张；一个域作为自身上的代数当然几何正则，但作为另一个底域上的代数就需要重新检查。第18章与第23章的不可分例子正说明底域不可省略。

\ProseParagraphBreak
卓越性的三个条件分别控制维数链、完备化纤维和正则点的邻域。完备局部环及域上有限型代数是主要来源；它们的卓越性与下面的稳定性质属于 Grothendieck 的卓越环理论。

\begin{theorem}\label{b3:thm:excellent-basic-classes}
完备 Noether 局部环是卓越环。若 \(A\) 卓越，则任意有限型 \(A\)-代数及其任意局部化也卓越；特别地，商环卓越。因而域上有限型代数及其局部化都是卓越环。
\end{theorem}
\begin{proofsketch}
完备 Noether 局部环的普适链性质采用松村英之的\cite[Theorem 29.4(ii), pp. 225--226]{Matsumura1989CommutativeRingTheory}。该证明先由 Cohen 结构定理把环表示为正则局部环的商，再用 Cohen--Macaulay 环的普适链定理；后一结果控制任意有限型代数中的素链长度，不能只用原环的一组参数的维数来代替。

\(G\) 环性质采用\cite[Theorem 32.3, pp. 257--258]{Matsumura1989CommutativeRingTheory}。这里需要检查每个局部化的形式纤维，而不只是恒等映射 \(A\to\widehat A=A\)。原证明先商去所考察纤维对应的素理想，化到完备局部整环，再用\cite[Theorem 29.4(iii), p. 226]{Matsumura1989CommutativeRingTheory}取完备正则局部子环 \(S\subseteq A\)，使 \(A\) 为有限 \(S\) 代数。令 \(K=\operatorname{Frac}S\)、\(L=\operatorname{Frac}A\)。经局部化、半局部完备化及直积分解，所考察的泛形式纤维是
\[
\bigl(\widehat{S_{\mathfrak q}}\otimes_{S_{\mathfrak q}}K\bigr)\otimes_KL
\]
的一个直积分量，其中 \(\mathfrak q\) 是所选素理想在 \(S\) 中的收缩。这说明为什么有限正则子环能够把问题归约到正则环的泛形式纤维。其几何正则性仍是关键步骤：正特征时须处理有限不可分扩张，原证明用相对形式光滑性及微分模的单射判据完成这一步。

完备局部环的 \(J\)-2 条件还采用半局部 Noether \(G\) 环满足 \(J\)-2 的定理；卓越性对有限型扩张、商及局部化的稳定性同样采用一般理论。松村在\cite[\S32, p. 260]{Matsumura1989CommutativeRingTheory}陈述这些结果，并指出多项式扩张保持 \(G\) 环性质的证明很深，转引其完整证明。上述形式纤维归约解释了这些定理中的一部分机制，正则点开放性和多项式扩张的比较定理则是这里所用而未展开的输入。最后，域是零维完备局部环，故域上有限型代数及其局部化的卓越性由前两项推出。
\end{proofsketch}

\subsection{Artin 逼近：域上有限型情形}
\label{b3:subsec:D04:02}

形式方程往往可以逐阶求解。若每一阶都能相容地继续，便得到完备化中的一个解；但这个解未必属于原环。逼近问题要求较少：给定一个有限阶数，能否在原环中找到一个准确满足方程的解，而且与这个形式解同余到该阶？

\begin{theorem}[Artin 逼近的代数情形]\label{b3:thm:artin-approximation-algebraic}
设 \(k\) 为域，\(R\) 为有限型 \(k\)-代数，\(\mathfrak p\subset R\) 为素理想，\(r,s\ge1\)。令 \((A,\mathfrak m)\) 为局部环 \(R_{\mathfrak p}\) 的 Hensel 化，\(\widehat A=\varprojlim_N A/\mathfrak m^N\) 为其 \(\mathfrak m\)-进完备化。
设
\[
F_1,\ldots,F_s\in A[Y_1,\ldots,Y_r],
\qquad
\widehat y\in\widehat A^{\,r},
\qquad F_i(\widehat y)=0\quad(1\le i\le s).
\]
则对每个整数 \(N\ge1\)，存在 \(y^{(N)}\in A^r\)，使
\[
F_i(y^{(N)})=0\quad(1\le i\le s),
\qquad
y^{(N)}-\widehat y\in\mathfrak m^N\widehat A^{\,r}.
\]
\end{theorem}

\begin{proofsketch}
本定理采用 Artin\cite[Theorem (1.10), p. 26]{Artin1969AlgebraicApproximation}的域上版本。原文把域的情形并入卓越 DVR 的情形，再于\cite[Section 5, pp. 42--44]{Artin1969AlgebraicApproximation}将一般 \(A\) 归约为多项式环在原点的 Hensel 化上的有限局部代数。有限展示把有限代数中的求解转成底环上的有限多项式方程组，这一步也处理原环中的零因子与幂零元。

证明的第一个主要输入是受控的光滑化归约。形式解定义同态 \(A[Y]/(F)\to\widehat A\)；原文在归约后的整环中取这个映射的关系理想，用卓越性保证的可分性取得泛光滑性，再用第4节的 N\'eron 型消解控制 DVR 参数处的行为\cite[Section 5, Lemma (5.6), pp. 44--46]{Artin1969AlgebraicApproximation}。没有关系时，任取足够近的元素即可；其余情形由此选出 \(G=(G_1,\ldots,G_h)\) 及一个 \(h\times h\) Jacobian 子式 \(\delta\)，使 \(\delta(\widehat y)\ne0\)。还须保留一个在形式解上非零的乘子 \(t\)，使 \(t\) 乘其余关系属于 \((G)\)；充分近的解保持 \(t(y)\ne0\)，于是整环中的 \(G(y)=0\) 推出全部关系\cite[Lemma (5.9), pp. 46--47]{Artin1969AlgebraicApproximation}。关系理想、可分性与光滑化的这些步骤采用原文的证明。

第二个主要输入是带整除条件的 Weierstrass 归纳\cite[Lemma (5.12), pp. 48--50]{Artin1969AlgebraicApproximation}。对 \(\delta(\widehat y)^2\) 作适当的坐标变换和 Weierstrass 预备，把未知级数分成次数有界的余式与商；余式系数满足少一个基底参数的有限多项式方程。先取 \(M\ge N\) 足够大，使同余于 \(\widehat y\) 模 \(\mathfrak m^M\) 的元素仍满足 \(t(y)\ne0\)；这样的 \(M\) 存在，因为 \(t(\widehat y)\ne0\)，而完备化是分离的。对余式方程归纳，并同时保留这一精度，得到 \(a\equiv\widehat y\pmod{\mathfrak m^M}\)，且
\[
G(a)\in\delta(a)^2\mathfrak m^M A^h.
\]
这里 \((A,\mathfrak m)\) 暂指归约后的 Hensel 局部整环，\(G\in A[Y_1,\ldots,Y_r]^h\)，其中 \(1\le h\le r\)。取得 \(\delta(a)^2\) 的整除条件是这一归纳的实质；只有误差属于很高的 \(\mathfrak m\) 次幂还不足以进行下面的修正。余式方程的构造及归纳的相容性采用所引引理。

说明最后的准确修正。对所选的 \(h\) 个未知量记 \(G\) 的 Jacobian 方阵为 \(J_G\)，\(\delta=\det J_G\)、\(J_G^{\mathrm{adj}}\) 为伴随矩阵。固定其余变量，代入
\(Y=a+\delta(a)J_G(a)^{\mathrm{adj}}U\)，其中这项修正只作用于所选 \(h\) 个坐标。由于 \(J_GJ_G^{\mathrm{adj}}=\delta I_h\)，Taylor 展开把 \(G(Y)=0\) 化为
\[
\delta(a)^2\bigl(e+U+Q(U)\bigr)=0,
\qquad e\in\mathfrak m^M A^h,
\]
其中 \(Q\in A[U_1,\ldots,U_h]^h\) 的各项关于 \(U\) 的次数至少二。对 \(e+U+Q(U)=0\)，零向量模 \(\mathfrak m\) 是简单零，Jacobian 为 \(I_h\)。由第23章的多变量 Jacobian 简单零提升\cref{b3:lem:henselian-jacobian-lifting}，取 \(I=\mathfrak m\)、初始点为零，可得 \(U\in\mathfrak mA^h\)；这也是原文所用的 Hensel 提升形式\cite[(1.9), pp. 25--26; Lemmas (5.10)--(5.11), pp. 47--48]{Artin1969AlgebraicApproximation}。又因 \(Q\) 没有常数项或一次项，可写 \(Q(U)=C(U)U\)，其中矩阵 \(C(U)\) 的各项属于 \(\mathfrak m\)。因此 \(I_h+C(U)\) 可逆，而
\[
U=-(I_h+C(U))^{-1}e\in\mathfrak m^MA^h.
\]
修正后的 \(y\) 便准确满足 \(G(y)=0\)，且 \(y\equiv\widehat y\pmod{\mathfrak m^M}\)，因而也同余到第 \(N\) 阶。所选精度保证 \(t(y)\ne0\)，故全部关系成立。最后经上述有限代数归约返回原环。这一伴随矩阵计算是原文 Lemma (5.10) 的修正机制，完整逼近定理还包含前面的光滑化与除法归纳。
\end{proofsketch}

更一般的卓越 Hensel 局部环上的 Artin 逼近还需一般的光滑化理论；上述原定理在这里采用的范围是域上有限型局部代数的 Hensel 化。

\ProseParagraphBreak
有两个量词必须保留。首先，\(y^{(N)}\) 对原方程是准确解，而非仅使误差属于 \(\mathfrak m^N\)。其次，\(N\) 改变时，允许选择不同的解；结论没有说存在一个 \(y\in A^r\) 与 \(\widehat y\) 在所有阶上同时相同。后一个要求在分离环中会直接迫使 \(y=\widehat y\)，强得多。

\begin{proposition}\label{b3:prop:linear-formal-approximation}
设 \((A,\mathfrak m)\) 为 Noether 局部环，\(\widehat A=\varprojlim_N A/\mathfrak m^N\) 为其 \(\mathfrak m\)-进完备化，
\(M\in\operatorname{Mat}_{s\times r}(A)\)、\(b\in A^s\)。若给定 \(\widehat y\in\widehat A^r\) 满足
\(M\widehat y=b\)，则对每个 \(N\ge1\)，存在准确解 \(y^{(N)}\in A^r\)，且
\(y^{(N)}-\widehat y\in\mathfrak m^N\widehat A^r\)。
\end{proposition}
\begin{proof}
取 \(a\in A^r\) 代表 \(\widehat y\) 模 \(\mathfrak m^N\) 的剩余类。令
\(L=M(\mathfrak m^NA^r)\subseteq A^s\)，这是有限生成模。完备化平坦且与有限生成模张量相容，所以
\[
L\widehat A=M(\mathfrak m^N\widehat A^r).
\]
误差 \(b-Ma=M(\widehat y-a)\) 属于这个子模，同时属于 \(A^s\)。由忠实平坦性，或有限生成模 \(A^s/L\) 到其完备化的单射性，可得 \(b-Ma\in L\)。故存在 \(c\in\mathfrak m^NA^r\)，使 \(Mc=b-Ma\)。取 \(y^{(N)}=a+c\)，便是所求准确解。
\end{proof}

线性情形只用有限生成模及完备化，不要求 Hensel 性。一般多项式方程中的变量乘积使误差不再由一个固定模映射控制；不能把以上证明直接套到非线性情形，也不能据此取消 Artin 定理对底环的假设。

\subsection{有限阶逼近、收敛性与代数化}
\label{b3:subsec:D04:03}

在 \(k[\![x]\!]\) 中，同余模 \((x)^N\) 表示前 \(N\) 个系数相同。它没有给剩余系数的大小任何限制，也没有断言一个任意形式关系能由多项式关系生成。

\begin{example}\label{b3:exm:approximation-does-not-converge}
在 \(\mathbb C[\![x]\!]\) 中取
\(h(x)=\sum_{j\ge1}j!x^j\)，令
\(\widehat y=(h^2,h)\)，则它准确满足 \(Y-Z^2=0\)。对每个 \(N\)，取多项式
\(h_N=\sum_{1\le j<N}j!x^j\)，便有准确的多项式解
\[
y^{(N)}=(h_N^2,h_N)
\equiv(h^2,h)\pmod{x^N}.
\]
这些准确解逼近一个发散形式解。不存在一个固定的收敛解与它在所有阶上相同，因为第二坐标的系数将被逐个确定为 \(j!\)。
\end{example}

代数化还需要说明保留什么结构。如果只问一个抽象完备环是否同构于某个代数局部环的完备化，形式坐标变换可以参与同构；若要求保留嵌入、指定函数、映射或标记，则必须把这些额外数据同时表示出来。\cref{b3:exm:formal-graph-not-analytic-equation}中的发散图像，恰好说明抽象环与指定嵌入可以有不同答案。

\ProseParagraphBreak
因此，有限阶数据最适合用来回答一个已经证明由有限阶决定的问题。例如已知一个商被 \((x)^N\) 零化时，该商的全部长度可由有限截断求出，见附录 F 的\secref{b3:sec:F04}；没有这个零化条件时，截断长度只是一个阶段性的量。有限决定性还需要额外的定理。

\subsection{奇点与消解理论导读}
\label{b3:subsec:D04:04}

结构定理使一个局部环成为正则环境中的方程商；标准基帮助计算其初始关系；形式纤维和逼近则比较不同局部对象。这些方法都能用于奇点，但它们本身不等于一个一般奇点消解定理。

\ProseParagraphBreak
最简单的平面曲线已经显示为什么要更换环境。对尖点
\(y^2=x^3\)，在原点外取新的比值坐标 \(v=y/x\)，代入 \(y=xv\) 得
\[
y^2-x^3=x^2(v^2-x).
\]
方程 \(v^2-x=0\) 定义光滑曲线，以 \(v\) 为参数，原坐标变成
\[
x=v^2,\qquad y=v^3.
\]
这与正规化的参数完全一致。保留全部乘积方程会额外留下由 \(x^2\) 定义的部分；取原点外曲线在新坐标图中的闭包，才得到严格变换 \(v^2-x=0\)。代数上，这是把总变换理想对 \(x\) 作饱和后收缩。

\ProseParagraphBreak
这一个坐标图是原点吹起的常见图。另一图取 \(x=yw\)，约去总变换中的 \(y^2\) 后得到 \(1-yw^3=0\)，在 \(y=0\) 上没有点，所以它不会增加尖点上方的新分支。两个图的相容关系解释了比值坐标如何合起来；若要把它作为一般几何构造，还须建立图的粘合、射影性及固有性等性质。

\ProseParagraphBreak
对于更高维对象，正规化通常不能消除全部奇点，局部图中的参数计算也不足以保证一种操作能在有限步后结束。一般消解理论要控制所选中心、例外部分和维数递降，并对特征及几何对象作准确限制。它所解决的是把奇异空间改造成正则空间的问题；自由消解则是在同一个底环上用自由模和正合复形描述模。二者都可参与奇点研究，却是不同的构造。
